\documentclass[11pt,a4paper]{article}
\usepackage[T1]{fontenc}
\usepackage[utf8]{inputenc}
\usepackage[english]{babel}
\usepackage{amsmath,amssymb}
\usepackage[margin=2cm]{geometry}
\setlength{\parindent}{0pt}
\setlength{\emergencystretch}{3em}
\allowdisplaybreaks

\title{About \texttt{false\_irrational\_equations.tex}}
\author{}
\date{}
\begin{document}
\maketitle
This note explains what the test data in \texttt{false\_irrational\_equations.tex} is used for. That file is input for the engine in \texttt{main.c}; it is read by the test \texttt{test\_irrational\_equations.c} in this folder.

\begin{enumerate}
\item \textbf{Purpose.} The file holds 30 false statements: the results of \texttt{book\_irrational\_equations.tex} with one detail changed.

\item \textbf{What it exercises.} The anti-regression check against false proofs: a solution dropped, a solution moved to a nearby number, or the extraneous root of the squared equation kept among the solutions. The engine must never prove one of them, and it finds a counterexample for each.

\item \textbf{Variety.} Three false variants for each of the ten equations. Each has a point where its two sides differ, found by an independent oracle (\texttt{tools/fuzz/fuzz\_radicals\_parameters.py}).

\item \textbf{Expected result.} The lines \texttt{PASS: 30 of 30 false statements rejected} and \texttt{PASS: 30 of 30 false statements refuted}.

\item \textbf{How to run it.} From the top folder of the repository:
\begin{verbatim}
cc -std=c11 -O2 -Wall -Wextra -o /tmp/test_irrational_equations \
    tests/equations_inequalities/test_irrational_equations.c \
    && /tmp/test_irrational_equations
\end{verbatim}
\end{enumerate}
\end{document}
